\chapter{Overview}\label{ch:overview}

This guide documents the competitive programming environment set up in the \file{Desktop/\allowbreak competitive} folder.
It is built around one idea: \textbf{the CPH extension stays in charge of test cases}. You fetch a problem, CPH creates
the file from your template and stores the samples, you write code, and you press CPH's \emph{Run All} button.
Everything else in this setup either makes that loop faster or explains what CPH cannot.

\section{The big picture}

\begin{figure}[H]
\centering
\begin{tikzpicture}[
  box/.style={draw=cpblue, fill=cpblue!6, rounded corners=3pt, thick, minimum width=3.2cm, minimum height=1.25cm,
    align=center, font=\sffamily\small},
  tool/.style={draw=cporange, fill=cporange!8, rounded corners=3pt, thick, minimum width=3.0cm, minimum height=1.05cm,
    align=center, font=\sffamily\small},
  arr/.style={->, thick, cpgray},
  lab/.style={font=\sffamily\scriptsize, text=cpgray, align=center}]
  \node[box] (browser) at (0, 0) {Problem page\\\scriptsize Codeforces, USACO, CSES};
  \node[box] (cc) at (4.5, 0) {Competitive\\Companion};
  \node[box, fill=cpteal!10, draw=cpteal] (cph) at (9, 0) {\textbf{CPH}\\\scriptsize stores the samples};
  \node[box] (file) at (9, -2.7) {\file{A\_Name.cpp}\\\scriptsize from \file{template.cpp}};
  \node[box] (editor) at (4.5, -2.7) {You write code\\\scriptsize templates, \texttt{dbg}};
  \node[tool] (f5) at (0, -2.7) {\key{F5}\ CPH tests\\\scriptsize crash reason, memory};
  \node[box, fill=cpgreen!10, draw=cpgreen] (run) at (4.5, -5.4) {CPH \emph{Run All}\\\scriptsize \key{Ctrl+Alt+B}};
  \node[box] (submit) at (9, -5.4) {Submit\\\scriptsize the file as it is};
  \node[tool] (stress) at (0, -5.4) {Stress test\\\scriptsize \key{Ctrl+Alt+X}};
  \draw[arr] (browser) -- node[above, lab] {click} (cc);
  \draw[arr] (cc) -- node[above, lab] {sends problem} (cph);
  \draw[arr] (cph) -- node[right, lab] {creates the file,\\fills the header\\and USACO I/O} (file);
  \draw[arr] (file) -- node[above, lab] {opens} (editor);
  \draw[arr] (editor) -- node[right, lab] {save} (run);
  \draw[arr] (run) -- node[above, lab] {all pass} (submit);
  \draw[arr, dashed] (editor) -- node[above=1pt, lab] {verdict is\\a mystery} (f5);
  \draw[arr, dashed] (f5) -- node[left, lab] {still stuck} (stress);
  \draw[arr, dashed] (stress) -- node[below=1pt, lab] {paste the\\failing input} (run);
\end{tikzpicture}
\caption{The daily loop. Solid arrows are the normal path; dashed arrows are what you do when a verdict surprises you.}
\end{figure}

\section{Quick start in 60 seconds}

\begin{enumerate}[label=\textbf{\color{cpteal}\arabic*.}]
  \item Open the \file{competitive} folder in VS Code. It uses the \emph{competitive} profile automatically.
  \item On a problem page, click the green Competitive Companion button. CPH creates \file{A\_Problem\_Name.cpp}
        from the template, fills in the problem name, link and time, and (for USACO) switches on file I/O.
  \item Write your solution in \cc{solve()}. Type \texttt{cp} to see every library snippet, or press
        \key{Ctrl+Alt+T} followed by a key (for example \key{U} for a DSU) to paste one at the cursor.
  \item Use \cc{dbg(x, v)} freely: it prints to CPH's \emph{Standard Error} box and vanishes on the judge.
  \item Press \key{Ctrl+Alt+B} (or CPH's \emph{Run All} button). If a verdict is confusing, press \key{F5}: the same tests
        run in the terminal with the crash reason, time and peak memory of each one.
  \item When all tests pass, submit the file exactly as it is. Nothing needs to be deleted or uncommented.
\end{enumerate}

\begin{tip}
The one-page summary of every shortcut is \file{CHEATSHEET.md} in the folder root. This guide is the long version;
the algorithms behind the templates are taught in the companion handbook, \emph{ALGORITHMS.pdf}.
\end{tip}

\section{What is where}

\newcommand{\desc}[1]{{\normalfont\sffamily\scriptsize\color{cpgray}--- #1}}
\begin{figure}[H]
\centering
\begin{forest}
  for tree={font=\ttfamily\small, grow'=0, child anchor=west, parent anchor=south, anchor=west, calign=first,
    edge path={\noexpand\path [draw, cpgray, \forestoption{edge}] (!u.south west) +(7.5pt,0) |- (.child anchor)\forestoption{edge label};},
    fit=band, s sep=1pt, before computing xy={l=15pt}}
  [{competitive/}
    [{template/template.cpp\quad\desc{the file CPH copies for every new problem}}]
    [{template/template.cpp.bak\quad\desc{your original template}}]
    [{include/debug.h\quad\desc{the dbg() printer (local builds only)}}]
    [{include/bits/\quad\desc{precompiled bits/stdc++.h (generated)}}]
    [{library/\quad\desc{129 templates in 10 category folders}}]
    [{tools/cp.py\quad\desc{test runner, stress tester, quickstart, setup, docs}}]
    [{stress/gen.cpp, brute.cpp\quad\desc{inputs for the stress tester}}]
    [{.vscode/tasks.json, settings.json\quad\desc{tasks behind F5 etc., IntelliSense}}]
    [{.cph/\quad\desc{CPH's test cases, one .prob file per problem (hidden)}}]
    [{docs/\quad\desc{this guide and the algorithms handbook (LaTeX + PDF)}}]
    [{archive/YYYY-MM/\quad\desc{old solutions, created by the archive shortcut}}]
    [{CHEATSHEET.md\quad\desc{one-page shortcut table}}]
  ]
\end{forest}
\caption{The folder layout. Build output (\texttt{.build/}, \texttt{*.exe}, \texttt{*.bin}) and \texttt{.cph/} are hidden from VS Code's explorer.}
\end{figure}

\section{What was changed outside the folder}

Keyboard shortcuts and snippets can only live in VS Code's user settings, so a few files outside the folder were
changed. All of them belong to the \emph{competitive} VS Code profile only; your default profile is untouched.

\begin{center}
\begin{tabularx}{\linewidth}{>{\raggedright\arraybackslash}p{5.2cm}L}
\toprule
\thd{File (in the profile folder)} & \thd{What changed} \\
\midrule
\file{keybindings.json} & One generated block, between \texttt{>>> cp-library shortcuts} and \texttt{<<<} markers, with the
  Ctrl+Alt+T / Ctrl+Alt+Y / Ctrl+Alt+U shortcuts, Ctrl+Alt+X and Ctrl+F5. Your own bindings above it are untouched.
  Backup: \file{keybindings.json.bak-cp-setup}. \\
\file{snippets/\allowbreak cp-library.code-snippets} & New file: the 175 snippets. \\
\file{settings.json} & Four CPH settings: compiler flags, template path, template variables on, stderr ignored for
  verdicts. Backup: \file{settings.json.bak-cp-setup}. \\
\midrule
TinyTeX & LaTeX packages were added (fonts, listings, forest, \ldots) to compile these PDFs. Additive only. \\
\bottomrule
\end{tabularx}
\end{center}

The profile folder is \file{\%APPDATA\%/Code/User/profiles/5d3986b6}. Re-running \emph{CP: Setup} rewrites only the
generated block and the four CPH settings, so it is always safe to run again.

\section{The four shortcuts that were broken before}

Your profile already had \key{F5}, \key{Ctrl+Shift+N}, \key{Ctrl+Shift+U} and \key{Ctrl+Shift+Alt+C} bound to tasks
named \texttt{cf-run}, \texttt{CP: Quickstart}, \texttt{CP: Quickstart USACO} and \texttt{CP: Archive Old Files},
but no such tasks existed anywhere, so the keys did nothing. They are now defined in \file{.vscode/tasks.json} and
work as described in Chapter~\ref{ch:testing} and Section~\ref{sec:quickstart}.

\section{How this setup is tied to CPH}\label{sec:cph-tie}

CPH is the centre of the workflow, so every piece was checked against CPH's own source code:

\begin{itemize}
  \item \textbf{CPH compiles with your flags.} CPH runs \texttt{g++ file.cpp -o file.bin \emph{Args} -D DEBUG -D CPH}.
        The \emph{Args} setting is written by \emph{CP: Setup} from one list in \file{tools/cp.py}, so the precompiled
        header, the 256\,MB stack, static linking and \cc{-DLOCAL} all apply to CPH runs.
        (The extra \texttt{-D DEBUG -D CPH} were verified not to disturb the precompiled header.)
  \item \textbf{CPH fills the template.} With template variables enabled, CPH replaces \texttt{\$name\$},
        \texttt{\$group\$}, \texttt{\$url\$}, \texttt{\$date\$}, \texttt{\$time\$}, \texttt{\$input\$} and
        \texttt{\$output\$} when it creates a file. That is how the header and the USACO file I/O are filled in.
  \item \textbf{CPH shows \cc{dbg} output.} CPH treats any stderr output as a failed test unless
        \texttt{ignoreSTDERROR} is on. It is now on, so \cc{dbg} output appears in the \emph{Standard Error}
        box while verdicts depend only on stdout and the exit code.
  \item \textbf{F5 reads CPH's tests.} The terminal runner reads the same \file{.cph/*.prob} file CPH uses.
  \item \textbf{The stress tester checks your brute force against CPH's samples} before it trusts it, and puts a
        failing input on the clipboard ready to paste into CPH as a new test case.
\end{itemize}

\begin{fact}[Honest limits]
Nothing is lost when you use only the CPH button. What CPH cannot do by design: tell you \emph{why} a program
crashed (it just fails the test), show memory use, or accept test cases from outside tools while its panel is open.
F5 covers the first two; the clipboard hand-off covers the third. The checked STL (\texttt{-D\_GLIBCXX\_DEBUG}) makes
local running times slower than the judge's, so treat CPH's times as upper bounds.
\end{fact}
